\begin{problem}[The nilradical sheaf]\label{prob:vi17-p16}
Let $N=\Nil(A)$ and let $X=\Spec A\cong\Spec(A/N)$ topologically.  Show that at each $\mfp$ the map on structure-sheaf stalks is
\[A_{\mfp}\to A_{\mfp}/N_{\mfp},\]
and explain how nilpotent sections are removed.
\end{problem}

\ifdefined\FullProblemDossiers
\paragraph{Why this problem matters.}
This identifies precisely how reduced and nonreduced structures can live on the same underlying space.

\paragraph{Useful ideas.}
Use localization commuting with quotients and the fact every prime contains $N$.

\paragraph{Guided hints.}
\begin{enumerate}
\item The corresponding prime of $A/N$ is $\mfp/N$.
\item Compute $(A/N)_{\mfp/N}$.
\end{enumerate}

% VI17_FULL_SOLUTION_REFINED
% VI17_FULL_SOLUTION_REFINED
\begin{solution}
Let $N=\Nil(A)$.  Since every prime ideal contains every nilpotent element, every prime $\mfp$ contains $N$.  Hence the quotient map $A\to A/N$ induces a bijection
\[
\Spec(A/N)\longrightarrow\Spec A,
\qquad
\mfp/N\longmapsto\mfp,
\]
which is in fact a homeomorphism.

At the corresponding point $\overline\mfp=\mfp/N$, the reduced structure sheaf has stalk
\[
(A/N)_{\overline\mfp}.
\]
Localization commutes with quotient, so there is a canonical isomorphism
\[
(A/N)_{\overline\mfp}\cong A_{\mfp}/N_{\mfp}.
\]
Therefore the morphism from the original structure sheaf to the reduced one induces on stalks exactly
\[
A_{\mfp}\longrightarrow A_{\mfp}/N_{\mfp}.
\]
Its kernel is the localized nilradical $N_{\mfp}$.

Consequently reduction does not alter the underlying point set or topology.  It changes the local function rings by killing precisely the nilpotent germs.  A nilpotent section may therefore be nonzero in the original structure sheaf while becoming zero everywhere in the reduced structure sheaf.
\end{solution}



\paragraph{What happened mathematically.}
Reduction changes functions, not points.

\paragraph{Extensions.}
Relate the kernels $N_{\mfp}$ to the sheaf of nilpotent sections and later to closed subschemes.

\paragraph{Related problems.}
VI.17.P15 and VI.17.P17 develop neighboring aspects of the same localization-sheaf calculus.

\paragraph{Provenance.}
\textbf{New synthesis of VI/10 and VI/17; no later closed-subscheme corpus Problem is consumed.}
\else
\paragraph{Provenance.} \textbf{New synthesis of VI/10 and VI/17; no later closed-subscheme corpus Problem is consumed.}
\fi
